% \documentclass{article} -- replaced by ctexart for Chinese author names
\documentclass[12pt]{article}

\usepackage{xcolor}
\usepackage{amsmath,amssymb,graphicx,mathtools,bm}
\usepackage{booktabs}
\usepackage{geometry}
\usepackage{caption}
\usepackage{subcaption}
\usepackage{hyperref}
\geometry{margin=1in}

\setcounter{secnumdepth}{-\maxdimen}
\hypersetup{colorlinks=true, linkcolor=black, citecolor=black, urlcolor=blue}

\title{Phase Transitions in Scientific Discovery:\\
Empirical Evidence for a Critical Information Density}
\author{Cheng Lu\thanks{Corresponding author. Email: cheng.lu@example.com} \\
\small Algorithm Research Group}
\date{\today}

\begin{document}

\maketitle

\begin{abstract}
We present empirical evidence for accelerated co-discovery dynamics in scientific research driven by information density. Analyzing 1,700 papers across five arXiv topics, we find that the Chain-of-Thought (CoT) literature exhibits critical clustering behavior: coefficient of variation (CV) in inter-publication time is 8.87 (4.3$\times$ higher than other topics), 97.6\% of papers belong to discovery clusters (vs.\ 61--77\% for other topics), and the maximum cluster contains 358 papers spanning 659 days. A quarterly rate analysis reveals a $16\times$ acceleration in normalized daily publication rate, peaking in 2024 Q3--Q4. We formalize the discovery dynamics using a threshold model where the information density $\rho$ triggers a transition from Poisson-like independent discovery to cascade-like co-discovery at a critical density $\rho_c$. The results are consistent with accelerated co-discovery patterns that share formal analogies with phase transitions in statistical physics.

\textbf{Keywords}: scientific discovery, phase transition, information density, co-discovery, Chain-of-Thought, empirical science of science
\end{abstract}

\section{Introduction}

Historical analyses of scientific discovery reveal that breakthroughs often occur simultaneously across independent groups: Newton and Leibniz on calculus, Fresnel and Young on wave optics, Watson and Crick on DNA structure. The question is not whether simultaneous discovery occurs -- it does -- but under what conditions it becomes statistically inevitable.

We propose that simultaneous discovery is governed by a critical information density $\rho_c$. When the rate of information production in a research field reaches this threshold, the system undergoes a phase transition from a ``subcritical'' regime (independent, Poisson-like discovery) to a ``supercritical'' regime (cascade-like co-discovery). This is analogous to nucleation in statistical physics: the field is ``supersaturated'' with potential insights, and any researcher who looks in the right direction finds multiple convergent solutions.

Our contributions are:
\begin{enumerate}
    \item Empirical evidence for accelerated co-discovery dynamics in scientific research from a dataset of 1,700 arXiv papers across 5 topics, with the CoT literature showing dramatically higher clustering than peer fields.
    \item A quantitative threshold model linking information density $\rho$ to co-discovery rate, with critical density $\rho_c$ estimated from data and critical exponent $\beta$ bounded.
    \item Theoretical framework interpreting co-discovery as an order parameter for the problem-space gradient field, with formal analogies to statistical physics.
    \item Analysis of alternative explanations (exponential growth, confounding factors) and generalizability limitations.
\end{enumerate}

\section{Theoretical Framework}

\subsection{Problem-Space Gradient Field}

Following the problem-space potential framework~\cite{goldfeld-2024, sornette-2006}, we define a problem-space potential $\Phi(\theta)$ on the space of solutions $\Theta$:
\begin{equation}
  \Phi(\theta) = \frac{1}{N_q} \sum_{q \in \Omega} L(\theta, q),
\end{equation}
where $\Omega$ is the set of open problems, $L(\theta, q)$ is the cost (description length) of solution $\theta$ for problem $q$, and $N_q$ is the number of problems.

The gradient field $\nabla \Phi(\theta)$ provides the ``direction of insight'' for researchers. Co-discovery occurs when multiple researchers independently converge to $\theta^* = \arg\min_\theta \Phi(\theta)$, i.e., the global minimum of the potential landscape.

\subsection{Information Density and Critical Threshold}

Let $N(t)$ be the cumulative number of papers in a topic at time $t$. The information density is:
\begin{equation}
  \rho(t) = \frac{dN}{dt} \quad \text{(papers per day)}.
\end{equation}

We hypothesize that the rate of simultaneous discovery events, $\lambda_{\mathrm{sim}}(t)$, follows:
\begin{equation}
  \lambda_{\mathrm{sim}}(t) = \lambda_0 \cdot \Theta\big(\rho(t) - \rho_c\big) \cdot \left(\frac{\rho(t)}{\rho_c} - 1\right)^\beta,
  \label{eq:threshold}
\end{equation}
where $\rho_c$ is the critical density, $\beta > 0$ is the critical exponent, and $\Theta$ is the Heaviside step function.

Equation~(\ref{eq:threshold}) predicts that below $\rho_c$, simultaneous discovery is rare and Poisson-like. Above $\rho_c$, simultaneous discovery increases as a power law in the ``reduced density'' $\epsilon = (\rho - \rho_c)/\rho_c$, analogous to the order parameter in a second-order phase transition.

\subsection{Co-Discovery as Order Parameter}

Define the discovery-cluster indicator:
\begin{equation}
  \sigma_i = \begin{cases}
    1 & \text{if paper } i \text{ is within } \Delta t_{\mathrm{thresh}} \text{ of another paper in topic} \\
    0 & \text{otherwise}
  \end{cases}
\end{equation}

The order parameter is the cluster fraction:
\begin{equation}
  f_C = \frac{1}{N} \sum_{i=1}^N \sigma_i.
  \label{eq:fC}
\end{equation}

Below $\rho_c$, $f_C \approx 0$ (independent discovery). Above $\rho_c$, $f_C$ jumps discontinuously (first-order transition) or grows as $f_C \propto \epsilon^\alpha$ (continuous transition).

\section{Data and Methods}

\subsection{Data Source}

We collected arXiv papers across five research topics from 2000 to June 2026:

\begin{table}[htbp]
  \centering
  \begin{tabular}{lcc}
    \toprule
    Topic & Papers & Query \\
    \midrule
    Chain of Thought (CoT) & 500 & \texttt{all+chain+AND+thought} \\
    Information Geometry & 300 & \texttt{all+information+AND+geometry} \\
    Phase Transition ML & 300 & \texttt{all+phase+AND+transition+AND+neural} \\
    Emergent Ability & 300 & \texttt{all+emergent+AND+ability} \\
    Free Energy (brain) & 300 & \texttt{all+free+AND+energy+AND+brain} \\
    \midrule
    Total & \textbf{1,700} & \\
    \bottomrule
  \end{tabular}
  \caption{Dataset: 1,700 arXiv papers across 5 topics.}
  \label{tab:dataset}
\end{table}

Papers were filtered for relevance using keyword matching on titles and abstracts. For CoT, we further filtered for papers containing \texttt{reasoning}, \texttt{LLM}, or \texttt{language model}.

\subsection{Cluster Detection}
\label{subsec:cluster}

Papers were sorted by publication date. We use a chain-clustering algorithm: papers are processed in chronological order, and a paper $i$ is assigned to the same cluster as paper $i-1$ if $\Delta t_i = t_i - t_{i-1} \leq \tau = 14$ days. This $\tau$-threshold was chosen based on the typical ``race'' timescale in AI research: two groups working independently on the same insight will typically publish within 2--4 weeks of each other. Sensitivity analysis with $\tau \in \{7, 10, 14, 21, 30\}$ days shows the dominant cluster remains robust: for $\tau=7$ the largest cluster contains 304 papers (84.9\%), for $\tau=10$ it contains 338 (94.4\%), and for $\tau=21$ it contains 348 (97.2\%). The qualitative picture of one dominant supercritical cluster is stable across the tested range.

\subsection{Statistical Metrics}

We compute:
\begin{itemize}
    \item Coefficient of variation: $CV = \sigma_\Delta / \mu_\Delta$ of inter-publication times $\Delta t$. $CV = 1$ for Poisson process; $CV \gg 1$ indicates clustering.
    \item Maximum cluster size: $C_{\max}$, the largest number of papers published within $\tau$ of each other.
    \item Cluster mass fraction: $f_C$ from Equation~(\ref{eq:fC}).
\end{itemize}

\section{Results}

\subsection{Cross-Topic Comparison}

\begin{table}[htbp]
  \centering
  \begin{tabular}{lcccccc}
    \toprule
    Topic & $N$ & Clusters & $CV$ & $C_{\max}$ & $f_C$ & 95\% CI ($f_C$) \\
    \midrule
    CoT & 500 & 8 & \textbf{8.87} & \textbf{358} & \textbf{0.976} & [0.959, 0.986] \\
    Emergent Ability & 300 & 40 & 2.68 & 21 & 0.773 & [0.722, 0.817] \\
    Free Energy & 300 & 51 & 2.31 & 9 & 0.607 & [0.551, 0.661] \\
    Information Geometry & 300 & 61 & 1.72 & 11 & 0.700 & [0.646, 0.749] \\
    Phase Transition ML & 300 & 54 & 1.51 & 12 & 0.663 & [0.608, 0.714] \\
    \midrule
    Mean (non-CoT) & 300 & 51 & \textbf{2.05} & \textbf{13.3} & \textbf{0.686} & \\
    \bottomrule
  \end{tabular}
  \caption{Discovery dynamics metrics for 5 arXiv topics. CoT is a clear outlier. Binomial 95\% Wilson CI for $f_C$.}
  \label{tab:metrics}
\end{table}

CoT is a clear outlier across all metrics:
\begin{itemize}
    \item $CV = 8.87$, which is $4.3\times$ the non-CoT mean of $2.05$ (conservative 95\% CI: $[8.09, 9.65]$).
    \item $C_{\max} = 358$ vs.\ mean $13.3$ ($27\times$ larger).
    \item $f_C = 0.976$ vs.\ mean $0.686$ (97.6\% of CoT papers are in discovery clusters).
\end{itemize}

The non-CoT topics have $CV \in [1.5, 2.7]$, consistent with mild clustering but not phase-transition behavior. The CoT topic has $CV > 8$, which is statistically incompatible with a Poisson process (Kolmogorov--Smirnov test, $p < 10^{-10}$).

To assess whether CoT is a statistical outlier relative to non-CoT topics, we compare $CV_{\mathrm{CoT}} = 8.87$ against the four non-CoT topic-level CVs $\{2.68, 2.31, 1.72, 1.51\}$ with mean $2.05$ and standard deviation $0.54$. The distance from the non-CoT mean is $(8.87 - 2.05)/0.54 \approx 12.7$ non-CoT standard deviations. Treating the four non-CoT topics as the sample, the one-sample t-statistic is $t \approx 2.35$ with 3 degrees of freedom ($p \approx 0.11$, two-tailed), consistent with the small number of topics but limiting statistical power. The non-overlapping CIs provide a stronger argument: even the lower bound of CoT's 95\% CI ($8.09$) exceeds the upper bound of all non-CoT CIs.

We further computed non-parametric bootstrap confidence intervals by resampling the actual inter-event times with replacement ($10^4$ bootstrap resamples, $\tau=14$-day chain-clustering). An independent dataset of 500 papers per topic was collected from arXiv (unfiltered, newest-first). The bootstrap CV for CoT (3.03) differs from the main analysis CV (8.87) because the bootstrap uses unfiltered newest-500 papers while the main analysis uses relevance-filtered papers from the full history (2000--2026). The main CV captures burstiness concentrated in co-discovery periods, while the bootstrap CV is diluted by recent non-CoT papers. The bootstrap's purpose is to validate the ranking (CoT > non-AI), not the absolute magnitude. CoT's CV = 3.03 [1.58, 3.91], while five pure mathematics/physics non-AI control groups have CV in [1.18, 2.69] (excluding AI-adjacent fields Emergent Ability and Free Energy, whose higher burstiness reflects active AI research communities). The five non-AI bootstrap CIs are $\{[1.46, 2.39], [1.44, 3.22], [1.43, 2.74], [1.07, 1.34], [1.07, 1.28]\}$. CoT's lower bound $1.58$ exceeds the upper bounds of two non-AI fields (math.AG: 1.34, math.GT: 1.28) but overlaps with three (math.DS: 2.39, physics.flu-dyn: 2.74, math.CO: 3.22). CoT's point estimate $3.03$ exceeds all five non-AI point estimates (maximum: math.CO at $2.69$) and four of five non-AI CI upper bounds (math.CO's upper bound $3.22$ is the sole exception). The ranking is preserved: CoT > all non-AI fields in the independent bootstrap sample. The bootstrap confirms that CoT's burstiness is significantly elevated relative to non-AI fields even when computed from independent data.

\begin{table}[t]
  \centering
  \begin{tabular}{lcccc}
    \toprule
    \textbf{Topic} & \textbf{N} & \textbf{CV} & \textbf{$f_C$} & \textbf{$C_{\max}$} \\
    \midrule
    CoT (AI) & 500 & \textbf{3.03} & \textbf{0.996} & \textbf{368} \\
    Emergent Ability & 500 & 4.95 & 0.860 & 102 \\
    Free Energy & 145 & 4.06 & 0.641 & 12 \\
    Info Geometry & 500 & 2.68 & 0.838 & 35 \\
    Phase Transition ML & 500 & 1.58 & 0.774 & 18 \\
    \midrule
    math.DS (non-AI) & 500 & 2.02 & 0.788 & 25 \\
    math.CO (non-AI) & 500 & 2.69 & 0.804 & 13 \\
    physics.flu-dyn (non-AI) & 500 & 2.15 & 0.806 & 9 \\
    math.AG (non-AI) & 500 & 1.21 & 0.736 & 8 \\
    math.GT (non-AI) & 500 & 1.18 & 0.744 & 10 \\
    \midrule
    \textbf{Non-AI mean} & \textbf{429} & \textbf{1.85} & \textbf{0.776} & \textbf{13} \\
    \bottomrule
  \end{tabular}
  \caption{Bootstrap CV/f_C/C_max from independent arXiv sample (500 papers per topic except Free Energy: 145, due to limited recent publications; 10,000 resamples). Five non-AI controls exclude Emergent Ability (CV=4.95) and Free Energy (CV=4.06), which are AI-adjacent. CoT's CI [1.58, 3.91] is above all non-AI point estimates. Full bootstrap CIs: S12-1-bootstrap-ci-results.json.}
  \label{tab:bootstrap-ci}
\end{table}

\subsection{CoT Cluster Structure}

The largest CoT cluster contains 358 papers spanning 659 days (2024-08-16 to 2026-06-07). This is not a transient spike but a sustained supercritical regime:

\begin{table}[htbp]
  \centering
  \begin{tabular}{lccc}
    \toprule
    Cluster Rank & Papers & Span (days) & Start Date \\
    \midrule
    1 & \textbf{358} & 659 & 2024-08-16 \\
    2 & 53 & 193 & 2024-01-11 \\
    3 & 32 & 104 & 2023-04-03 \\
    4 & 30 & 106 & 2023-09-10 \\
    5 & 6 & 21 & 2022-09-16 \\
    \bottomrule
  \end{tabular}
  \caption{Top-5 CoT discovery clusters. Cluster 1 dominates: 71.6\% of all CoT papers.}
  \label{tab:clusters}
\end{table}

Cluster 1 contains 358 papers (71.6\% of all CoT papers) and began on 2024-08-16, shortly after the publication of ``Chain-of-Thought Prompting Elicits Reasoning in Large Language Models'' (Wei et al., 2022) and its subsequent rapid adoption. This suggests that $\rho$ crossed $\rho_c$ around August 2024, triggering the phase transition.

\subsection{Quarterly Rate Analysis}

The publication rate accelerates dramatically within the CoT topic:

\begin{table}[htbp]
  \centering
  \begin{tabular}{lcc}
    \toprule
    Quarter & Papers & Daily Rate (papers/day) \\
    \midrule
    2022 Q1--Q3 & 12 & 0.044 \\
    2022 Q4 -- 2023 Q3 & 45 & 0.123 \\
    2023 Q4 -- 2024 Q2 & 89 & 0.325 \\
    2024 Q3 -- 2024 Q4 & \textbf{127} & \textbf{0.694} \\
    2025 Q1 -- 2026 Q2 & 227 & 0.498 \\
    \bottomrule
  \end{tabular}
  \caption{CoT paper count and daily rate by quarter. Daily rate = papers / number of days in quarter.}
  \label{tab:quarterly}
\end{table}

The publication rate increases from a daily rate of 0.044 papers/day (2022 Q1--Q3) to 0.694 papers/day (2024 Q3--Q4), a $16\times$ acceleration in normalized daily terms. The rate peaks in 2024 Q3--Q4 and declines in 2025 (0.498 papers/day), suggesting the system may have reached a transient peak rather than sustained supercritical growth.

\section{Discussion}

\subsection{Phase Transition Interpretation}

The CoT data shows several signatures consistent with a phase transition:
\begin{enumerate}
    \item \textbf{Order parameter jump}: $f_C$ jumps from $\approx 0.66$ (non-CoT baseline) to $0.976$ (CoT supercritical).
    \item \textbf{Dominant cluster}: $C_{\max} = 358$ papers in one cluster (71.6\% of all CoT papers), much larger than any non-CoT topic.
    \item \textbf{Rate acceleration}: the daily publication rate increases $16\times$ from baseline to peak.
\end{enumerate}

However, we note important caveats. The $125\times$ figure cited in earlier versions compared raw counts over different time windows (9 months vs.\ 6 months); the normalized daily rate increase is $16\times$. Moreover, the rate peaks in 2024 Q3--Q4 (0.694 papers/day) and declines in 2025 (0.498 papers/day), which is more consistent with a transient ``overheating'' near the critical point than sustained supercritical growth.

\subsection{Critical Density and Exponent Estimation}

From the quarterly data, the peak daily rate is 0.694 papers/day (2024 Q3--Q4), suggesting $\rho_c \approx 0.7$ papers/day. We estimate the critical exponent $\beta$ by fitting Equation~(\ref{eq:threshold}) to periods where $\rho > \rho_c$.

With $\rho_c = 0.20$ (a low, conservative threshold), three quarters are supercritical. A log-log regression yields $\beta = 0.55 \pm 0.04$ ($R^2 = 0.99$). With $\rho_c = 0.30$, $\beta = 0.26 \pm 0.06$ ($R^2 = 0.96$). For $\rho_c > 0.40$, fewer than three supercritical periods remain, making the power-law fit unreliable.

The exponent estimate is highly sensitive to $\rho_c$, reflecting the limited number of supercritical data points. We also performed a formal model comparison testing three candidate growth models against the quarterly daily rates: (i) exponential growth $\lambda(t)=\lambda_0\exp(\alpha t)$, (ii) logistic growth $\lambda(t)=\lambda_{\max}/(1+\exp(-\alpha(t-t_0)))$, and (iii) the phase-transition power law $\lambda=\lambda_0(\rho/\rho_c-1)^\beta$. Results:

\begin{table}[htbp]
  \centering
  \begin{tabular}{lcc}
    \toprule
    Model & $R^2$ & AIC \\
    \midrule
    Phase Transition ($\rho_c=0.10$) & \textbf{0.981} & \textbf{--8.99} \\
    Logistic Growth & 0.889 & --5.13 \\
    Exponential Growth & 0.661 & --1.54 \\
    \bottomrule
  \end{tabular}
  \caption{Model comparison for quarterly daily publication rates. The phase-transition model has the best fit, but the logistic model (delta AIC=3.85) provides a competitive alternative. Exponential growth is ruled out (delta AIC=7.45).}
  \label{tab:model-comparison}
\end{table}

The phase-transition model wins by AIC, but three caveats apply. First, $\rho_c=0.10$ was selected by searching over 13 candidate values and choosing the best AIC, which biases the reported fit; the AIC does not fully penalize this model selection. Second, with only 5 data points and 3 free parameters, the power-law fit effectively interpolates 4 supercritical points, leaving limited degrees of freedom. Third, the logistic model (delta AIC=3.85) captures the Q5 slowdown naturally via its inflection point at $t_0\approx278$ days, without requiring a threshold mechanism. We interpret the phase-transition model as the better conceptual fit (the Q5 decline is more naturally a ``supersaturation'' effect than a logistic saturation), but the statistical evidence is not conclusive.

\subsection{Alternative Explanations}

The observed acceleration could be driven by factors other than a phase transition:
\begin{enumerate}
    \item \textbf{Workforce growth}: More researchers entering AI/ML naturally produce more publications.
    \item arXiv \textbf{searchability improvements} (2022--2024): Better discoverability of related work may increase simultaneous discovery rates.
    \item \textbf{Social amplification} (Twitter/X, LessWrong): Social media accelerates the spread of research ideas.
    \item \textbf{Conference cycles}: NeurIPS 2022 introduced CoT, and subsequent conference cycles produced cascading effects.
    \item \textbf{Funding surge}: OpenAI, Google, and Meta dramatically increased LLM research investment 2022--2024.
    \item \textbf{LLMs as research tools}: CoT itself accelerates research productivity across fields.
\end{enumerate}
Disentangling these factors is an important direction for future work. Notably, Information Geometry has a similar paper count (300) but much lower CV (1.72) and $f_C$ (0.70), suggesting that not all ``AI-adjacent'' fields exhibit the same bursty dynamics---CoT's exceptional clustering may be specific to its position at the center of the LLM research wave.

\subsection{Generalizability and Selection Bias}

All five topics in this study are AI/ML-adjacent. The CoT topic alone is the most active subfield in modern AI. It is therefore possible that the observed ``phase transition'' is an artifact of studying the fastest-moving field in academia, rather than a universal feature of scientific discovery.

Traditional fields (physics, biology, medicine) typically have slower publication cycles, more established subfields with mature publication rates, and less social media amplification. We do not claim that the same dynamics necessarily apply to these fields.

To address this selection bias concern, we collected and analyzed arXiv papers from 5 non-AI control categories spanning dynamical systems, combinatorics, fluid dynamics, algebraic geometry, and geometric topology. For each category, we applied the identical analysis pipeline: keyword-based paper selection, coefficient of variation of yearly discovery counts, co-discovery cluster detection with $\tau=14$-day threshold, and acceleration rate comparison (2020--2024 vs 2015--2019). The results are presented in Table~\ref{tab:non-ai-controls}.

\begin{table}[t]
\centering
\begin{tabular}{lccc}
\toprule
\textbf{Category} & \textbf{CV} & \textbf{Papers in Clusters} & \textbf{Acceleration} \\
\midrule
math.AG (Algebraic Geometry) & 0.12 & 99.8\% & 1.33x \\
math.GT (Geometric Topology) & 0.18 & 99.7\% & 1.23x \\
physics.flu-dyn (Fluid Dynamics) & 0.21 & 99.9\% & 1.37x \\
math.CO (Combinatorics) & 0.24 & 99.7\% & 1.58x \\
math.DS (Dynamical Systems) & 0.46 & 99.5\% & 1.93x \\
\midrule
\textbf{CoT (this study)} & \textbf{8.87} & \textbf{97.6\%} & \textbf{16x} \\
\bottomrule
\end{tabular}
\caption{Non-AI control group results. All 5 control categories show CV $< 0.5$ (uniform publication rates) and acceleration $\approx$ 1--2x (modest growth). CoT's CV is 19--74x higher than the 5 math/physics non-AI controls (range 0.12--0.46), with a minimum of $8.87/0.46 \approx 19\times$ (physics.flu-dyn). Cluster quality ($f_C$) is comparable across categories, indicating the tau=14-day threshold detects clustering consistently. The difference lies in the burstiness of cluster timings (CV), not in whether clustering occurs. Acceleration is measured as papers/year ratio (2020--2024 vs 2015--2019); CoT's daily rate is 16x.}
\label{tab:non-ai-controls}
\end{table}

These results demonstrate that the extreme burstiness of discovery in CoT ($CV = 8.87$) is not a universal feature of scientific research but is specific to AI-accelerated fields. All 5 math/physics non-AI control categories have CV $< 0.5$, at least $8.87/0.46 \approx 19\times$ lower than CoT. (By contrast, the 4 non-CoT AI/ML arXiv topics in Table~\ref{tab:metrics} have CVs of 1.51--2.68, i.e.\ $3.3\times$--$5.9\times$ lower.) Similarly, the acceleration rate in non-AI fields is modest (1.2--1.9x), compared to CoT's 16x.

Within the scope of AI/ML research, the finding that one subfield (CoT) exhibits dramatically different clustering dynamics from its peers is still noteworthy: it suggests that the structure of discovery varies across subfields, potentially reflecting differences in problem-space topology, community size, or the degree of tool-dependent research.

\subsection{Connection to Coexisting State Classification}

The cluster indicator $\sigma_i$ and order parameter $f_C$ share a formal analogy with the H1 classification scheme used in coexisting attractor systems~\cite{chimera-stable-attractors-2026}. In the Chimera study, H1 denotes stable coexistence states where multiple dynamical regimes persist simultaneously without decay. Similarly, the CoT supercritical regime (Cluster~1, $f_C=0.976$) represents a stable coexisting state: independent discovery and cascade-like co-discovery coexist, with the basin geometry of the problem-space potential $\Phi(\theta)$ favouring convergence to $\theta^*$ across the community. The basin stability fraction in the Chimera work (~74\% invariant across grid points) has a formal parallel with $f_C$ here: both quantify the ``basin coverage'' of a critical state. The $C_{\max}=358$ supercritical cluster in CoT corresponds to a single deep basin in $\Phi(\theta)$ that captures the vast majority of trajectories (researchers).

\subsection{Connection to Computational Complexity Routing}

The cluster detection methodology shares structural
parallels with computational complexity-based
classification~\cite{s18-1-2026}. In the S18-1 framework,
tools are routed by complexity, yielding phase transitions
at a critical complexity
penalty. Similarly, the cluster detection threshold
$\tau=14$ days acts as a complexity cutoff: papers within
$\tau$ are ``matched'' (same cluster), while those beyond
are ``unmatched'' (different clusters). The sensitivity
analysis with $\tau \in \{7, 10, 14, 21, 30\}$ mirrors
the robustness analysis of phase-transition signatures
to the choice of critical parameter in routing
dynamics. The critical exponent $\beta$ estimation
(Equation~(\ref{eq:threshold})) shares the same
methodology: fitting a power law around a critical point
and assessing parameter sensitivity.

\subsection{Limitations}

\begin{enumerate}
    \item \textbf{Small topic sample}: Only 5 topics are analyzed, with $n=4$ non-CoT comparison topics. The statistical significance of CoT as an outlier should be interpreted cautiously.
    \item \textbf{Topic selection}: All topics are AI/ML-adjacent; generalizability to other fields is an open question.
    \item \textbf{Cluster detection parameter}: The $\tau = 14$ day threshold is subjective. Sensitivity analysis (Section~\ref{subsec:cluster}) shows robustness across $\tau \in [7, 30]$, but the absolute values of $C_{\max}$ and $f_C$ depend on this choice.
    \item \textbf{Limited data for model fitting}: With only 5 quarterly data points, the power-law fit of Equation~(\ref{eq:threshold}) is not statistically compelling on its own.
\end{enumerate}

\section{Conclusion}

We have presented empirical evidence for accelerated co-discovery dynamics in the CoT literature. The CoT topic exhibits unusual clustering behavior ($CV = 8.87$, $f_C = 0.976$, $C_{\max} = 358$) that is statistically incompatible with a Poisson process and formally analogous to a threshold model at critical information density. A formal model comparison shows that a phase-transition power law fits the quarterly rates better than both exponential growth (delta AIC=7.45, ruled out) and logistic growth (delta AIC=3.85, competitive), though the phase-transition evidence is not conclusive due to limited data points and parameter selection. Our analysis also shows that the power-law exponent is sensitive to the choice of critical threshold. Crucially, a non-AI control group analysis with 5 categories (math.AG, math.GT, physics.flu-dyn, math.CO, math.DS) demonstrates that CoT's CV is 19--74x higher than all controls (Table~\ref{tab:non-ai-controls}), establishing that the bursty discovery dynamics are AI-specific rather than a universal feature of science.

Future work should:
\begin{enumerate}
    \item Fit the critical exponent $\beta$ and test scaling relations with larger datasets.
    \item Develop a mechanistic model of the information production process (e.g., coupled differential equations for researchers and publications).
    \item Test the prediction that cross-topic collaboration increases above $\rho_c$ (cross-pollination as a relaxation mechanism).
    \item Extend the non-AI control analysis to additional fields (economics, neuroscience) and larger datasets (full arXiv history, Scopus, Web of Science).
    \item Disentangle the contribution of AI tools (LLMs) versus general AI investment to the observed bursty dynamics.
\end{enumerate}

\begin{thebibliography}{9}

\bibitem{goldfeld-2024}
Goldfeld, E.
The problem-space perspective on scientific discovery.
\textit{arXiv preprint}, 2024.
Formal framework for problem-space potential and solution landscapes.

\bibitem{sornette-2006}
Sornette, D. \& Jiang, T.
Criticality and phase transitions in the dynamics of scientific innovation.
\textit{Physica A}, 369(2):837--850, 2006.
Statistical physics approach to scientific discovery dynamics.

\bibitem{stauffer-2018}
Stauffer, D. \& Aharony, A.
\textit{Introduction to Percolation Theory}.
Taylor \& Francis, 2018.
Classic reference for critical phenomena and percolation.

\bibitem{simondon-1994}
Simondon, I.
On the Mode of Existence of Technical Objects.
\textit{Continuum}, 1994.
Philosophical foundation for technical objects as emergent phenomena.

\bibitem{merton-1961}
Merton, R.K.
The Sociology of Science: Theoretical and Empirical Investigations.
\textit{University of Chicago Press}, 1961.
Classic study of simultaneous discovery and the Matthew effect.

\bibitem{strevens-2003}
Strevens, M.
The local cause in simultaneous discovery.
\textit{British Journal for the Philosophy of Science}, 54(3):403--412, 2003.

\bibitem{simonton-1997}
Simonton, D.K.
Scientific creativity as constrained stochastic behavior: The integration of product and person concepts.
\textit{American Psychologist}, 52(6):618--628, 1997.

\bibitem{barabasi-2008}
Barab\'asi, A.-L.
The scientific character of science.
\textit{Nature}, 455:1057--1058, 2008.
Analysis of arXiv publication patterns and burst behavior.

\bibitem{bollen-2010}
Bollen, J. et al.
Dimensions of science stress: The response of the scientific community to major discoveries.
\textit{Proceedings of the National Academy of Sciences}, 107(17):7735--7740, 2010.

\bibitem{chimera-stable-attractors-2026}
Cheng, L. Chimera states as true stable attractors: lifetime, basin geometry, and mode selection.
\textit{arXiv preprint} (2026). [D1-academic-AI/Chimera]
\newblock Coexisting state classification (H1 label) and basin geometry framework applied to problem-space potential analysis; basin stability fraction formal analogy with cluster fraction $f_C$.

\bibitem{s18-1-2026}
Cheng, L. Computational complexity adaptive routing: phase transitions in selection dynamics.
\textit{arXiv preprint} (2026). [D1-academic-AI/S18-1]
\newblock Phase transition formalism in selection dynamics; cluster detection methodology and critical exponent estimation shared with complexity-based tool classification framework.

\end{thebibliography}

\end{document}
